\documentclass[10pt]{article}
\usepackage[T1]{fontenc}
\usepackage[utf8]{inputenc}
\usepackage[french]{babel}
\usepackage{amsmath,amssymb,mathrsfs,mathpazo}
\usepackage{geometry}
\geometry{a4paper,margin=2cm}
\usepackage{graphicx,array,tabularx,tabularray,colortbl,multirow,multicol,booktabs,enumitem,pifont,ulem,fancybox,diagbox,cancel}
\usepackage{pstricks,pstricks-add,pst-plot,pst-tree,pst-node,pst-text,pst-func,pst-math,pst-bspline,pst-3dplot,pst-eucl,pst-solides3d}
\let\pscalculate\undefined
\usepackage{pst-bezier}
\usepackage{tikz,pgfplots,tkz-tab}
\usetikzlibrary{arrows,arrows.meta,patterns,calc,positioning,trees,decorations.markings,shapes.geometric,shapes.misc,shapes.arrows,intersections,decorations.pathmorphing,decorations.pathreplacing,backgrounds,shadows,babel}
\pgfplotsset{compat=1.18}
\usepgfplotslibrary{fillbetween}
\usepackage[np]{numprint}
\usepackage{listings,makecell,pas-tableur,twoopt}
\let\barre\undefined
\usepackage[pstricks]{tablvar}
\definecolor{darkgreen}{RGB}{0,100,0}
\providecommand{\red}{\color{red}}
\providecommand{\blue}{\color{blue}}
\providecommand{\green}{\color{green}}
\providecommand{\black}{\color{black}}
\providecommand{\white}{\color{white}}
\providecommand{\gray}{\color{gray}}
\providecommand{\cyan}{\color{cyan}}
\providecommand{\magenta}{\color{magenta}}
\providecommand{\yellow}{\color{yellow}}
\providecommand{\violet}{\color{violet}}
\providecommand{\orange}{\color{orange}}
\providecommand{\euro}{\texteuro}
\providecommand{\EUR}{\texteuro}
\providecommand{\ieme}{\textsuperscript{e}}
\providecommand{\degres}{\ensuremath{^\circ}}
\providecommand{\degree}{\ensuremath{^\circ}}
\providecommand{\vv}[1]{\overrightarrow{#1}}
\providecommand{\Coord}[1]{#1}
\providecommand{\decosix}{}
\providecommand{\starredbullet}{\ensuremath{\bullet}}
\pagestyle{empty}
\setlength{\parindent}{0pt}
\usepackage{adjustbox}\DeclareUnicodeCharacter{2236}{:}\DeclareUnicodeCharacter{2190}{\ensuremath{\leftarrow}}
\providecommand{\R}{\mathbb{R}}
\providecommand{\N}{\mathbb{N}}
\providecommand{\Z}{\mathbb{Z}}
\providecommand{\Q}{\mathbb{Q}}
\providecommand{\D}{\mathbb{D}}
\providecommand{\C}{\mathbb{C}}
\providecommand{\e}{{\mathrm{e}}}
\providecommand{\vect}[1]{\overrightarrow{#1}}
\providecommand{\vectt}[1]{\overrightarrow{\mathrm{#1}}}
\providecommand{\vvt}[1]{\overrightarrow{\mathrm{#1}}}
\providecommand{\np}[2][]{{#2~#1}}
\providecommand{\vv}[1]{{\overrightarrow{#1}}}
\renewcommand{\d}{{\mathrm{d}}}
\providecommand{\degres}{{^\circ}}
\providecommand{\degre}{{^\circ}}
\providecommand{\blue}{\color{blue}}
\providecommand{\red}{\color{red}}
\providecommand{\enskip}{\quad}
\providecommand{\llbracket}{\lbrack\!\lbrack}
\providecommand{\rrbracket}{\rbrack\!\rbrack}
\providecommand{\euro}{{\text{€}}}
\providecommand{\no}{{\text{n°}}}
\providecommand{\interval}[2]{{[#1;#2]}}
\providecommand{\intervFO}[2]{{[#1;#2[}}
\providecommand{\dotfill}{\ldots}

\providecommand{\ds}{\displaystyle}
\providecommand{\textding}[1]{#1}
\providecommand{\Oij}{\(\mathrm{O};\vec{\imath},\vec{\jmath}\)}
\providecommand{\Oijk}{\(\mathrm{O};\vec{\imath},\vec{\jmath},\vec{k}\)}
\providecommand{\Ouv}{\(\mathrm{O};\vec{u},\vec{v}\)}
\providecommand{\intervFF}[2]{\ensuremath{[#1;#2]}}
\renewcommand{\euro}{\ifmmode\text{\texteuro}\else\texteuro\fi}
\DeclareUnicodeCharacter{00B0}{\ifmmode^{\circ}\else\textdegree\fi}
\DeclareUnicodeCharacter{20AC}{\ifmmode\text{\texteuro}\else\texteuro\fi}
% Énoncés : Éducation nationale. Transcription et corrigés : APMEP, Denis Vergès et contributeurs. Découpage et classement : CoopMaths, Sébastien Lozano, Éric Elter et contributeurs. Adaptation de la présentation : Maths Explorer.
% CC BY-SA : https://coopmaths.fr/www/about/
% https://coopmaths.fr/alea/static/e3c/2020/tex/e3c_2020_00_sujet13_3.tex
\begin{document}
\noindent\textit{APMEP / CoopMaths. CC BY-SA. Présentation adaptée par Maths Explorer.}\par\medskip
\medskip
Un cafetier propose à ses clients des cookies au chocolat ou aux noisettes en s'approvisionnant dans trois boulangeries. Un client prend un cookie au hasard.

On note:

$C$ l'évènement \og le cookie est au chocolat \fg,

$N$ l'évènement \og le cookie est aux noisettes \fg,

$B_1$ l'évènement \og le cookie provient de la boulangerie 1 \fg, 

$B_2$ l'évènement \og le cookie provient de la boulangerie 2 \fg,

$B_3$ l'évènement \og le cookie provient de la boulangerie 3 \fg.

On suppose que :

\setlength\parindent{1cm}
\begin{itemize}
\item la probabilité que le cookie provienne de la boulangerie 1 est de 0,49 ;
\item la probabilité que le cookie provienne de la boulangerie 2 est de 0,36 ;
\item  $P_{B_2}(C) = 0,4$ où $P_{B_2}(C)$ est la probabilité conditionnelle de $C$ sachant $B_2$ ;
\item la probabilité que le cookie soit aux noisettes sachant qu'il provient de la troisième boulangerie est de $0,3$.
\end{itemize}
\setlength\parindent{0cm}

L'arbre pondéré ci-dessous correspond à la situation et donne une information
supplémentaire : le nombre $0,6$ sur la branche de $B_1$ à $C$.

\begin{center}
\pstree[treemode=R,nodesepA=0pt,nodesepB=3pt]{\TR{}}
{\pstree{\TR{$B_1$~~}\naput{\ldots}}
	{\TR{$C$} \naput{0,6}
	\TR{$N$}\tbput{\ldots}
	}
\pstree{\TR{$B_2$~~}\naput{\ldots}}
	{\TR{$C$} \taput{\ldots}
	\TR{$N$}\tbput{\ldots}
	}
\pstree{\TR{$B_3$~~}\nbput{\ldots}}
	{\TR{$C$}\taput{\ldots} 
	\TR{$N$}\tbput{\ldots}
	}	
}
\end{center}

\smallskip

\begin{enumerate}
\item Exprimer par une phrase l'information donnée par le nombre 0,6 sur la branche de $B_1$ à $C$.
\item Recopier et compléter sur la copie l'arbre pondéré ci-dessus.
\item Définir par une phrase l'évènement $B_1 \cap C$ et calculer sa probabilité.
\item Montrer la probabilité $P(C)$ d'avoir un cookie au chocolat est égale à $0,543$.
\item Calculer la probabilité d'avoir un cookie provenant de la boulangerie 2 sachant qu'il est au chocolat. On donnera le résultat arrondi au millième.
\end{enumerate}

\bigskip

\end{document}
